% Modernised reading — fonds Alexandre Grothendieck, shelfmark 108, pages 1-5.
%
% This is an INTERPRETATION, not a transcription. Everything the critical
% apparatus carried moves into footnotes. In any dispute of substance the
% transcription governs: whoever cites Grothendieck must cite batch-01.fr.tex.
%
% Pass: Opus 5 (claude-opus-5), 2026-08-16 — first pass, unchecked against the
% pages by a human. The folder was read whole; it holds one batch, pages 1-5.
%
% The standing risk in this folder is priority. Several of its seventeen
% questions have since been answered, some of them famously, and it would be easy
% to write a reading in which Grothendieck is seen anticipating each answer. That
% is not done here. A question is stated in modern terms and named; whether it
% was settled, and by whom, is said only where saying it costs nothing — and
% never in a form that implies he knew.
\documentclass[11pt,a4paper]{article}
% Le préambule commun à toutes les transcriptions du fonds Grothendieck.
%
% Il tient en une idée : **l'incertitude se compose, elle ne se résout pas.**
% Une transcription qui lisse un mot douteux a détruit la seule chose pour
% laquelle on la faisait — un lecteur doit pouvoir savoir, à chaque endroit,
% ce qui était sur le feuillet et ce qui a été ajouté. Les six macros qui
% suivent sont donc l'appareil critique, et non de la décoration.
%
% Les mêmes commandes sont reconnues par `scripts/render.mjs`, qui produit la
% vue de lecture du site. Une macro ajoutée ici doit l'être aussi là-bas,
% faute de quoi le rendu échouera bruyamment — ce qui est voulu : un rendu
% silencieusement faux est pire qu'un rendu absent.

\ProvidesPackage{grothendieck}[2026/08/08 Transcriptions du fonds Grothendieck]

\usepackage{amsmath,amssymb,amsthm}
\usepackage{mathrsfs}
\usepackage{stmaryrd}
% Les diagrammes commutatifs sont partout dans ce fonds : c'est la langue
% même de la géométrie algébrique de Grothendieck, et une transcription qui
% les rendrait en prose ne transcrirait rien.
\usepackage{tikz-cd}
% Français par défaut — c'est la langue des feuillets — mais l'anglais est
% chargé aussi : la traduction et le résumé sont des documents anglais, et la
% typographie française (espaces avant les deux-points) y serait une faute.
% Ces documents basculent par \selectlanguage{english} en tête de corps.
\usepackage[english,french]{babel}
\usepackage{fontspec}
\usepackage{geometry}
\geometry{a4paper,margin=25mm,marginparwidth=28mm}
\usepackage{marginnote}
\usepackage{xcolor}
% ulem, not soul. `soul`'s \st and \ul are built on a character-by-character
% scan that XeLaTeX's font machinery breaks: the document fails with an
% undefined control sequence rather than degrading. ulem does the same two jobs
% and is Unicode-engine safe. `normalem` stops it hijacking \emph, which the
% transcriptions use for Grothendieck's own emphasis.
\usepackage[normalem]{ulem}
\usepackage{hyperref}
\hypersetup{colorlinks=true,linkcolor=black,urlcolor=black,pdfborder={0 0 0}}

% --- Métadonnées du lot -----------------------------------------------------
% Reprises telles quelles par le script de rendu, qui les affiche en tête de
% la vue de lecture. Elles fixent ce que le fichier prétend couvrir : sans
% elles, une transcription flotte sans point d'attache dans un fonds de seize
% mille feuillets.

\newcommand{\folder}[1]{\gdef\grothFolder{#1}}
\newcommand{\batch}[1]{\gdef\grothBatch{#1}}
\newcommand{\leaves}[2]{\gdef\grothFirst{#1}\gdef\grothLast{#2}}
\newcommand{\foldertitle}[1]{\gdef\grothFoldertitle{#1}}
\gdef\grothFolder{?}\gdef\grothBatch{?}\gdef\grothFirst{?}\gdef\grothLast{?}\gdef\grothFoldertitle{}

% --- Le résumé --------------------------------------------------------------
%
% Il ouvre la lecture modernisée, sous le titre « Résumé », et il est composé
% en petit. Ce n'est pas un ornement : c'est le seul passage du document qui
% ne soit pas des mathématiques, et un lecteur doit pouvoir le reconnaître
% avant de l'avoir lu — puis le sauter s'il connaît déjà le sujet, ou s'y
% arrêter s'il ne le connaît pas. Le filet qui le suit marque la reprise du
% corps.

\newenvironment{resume}
  {\small\setlength{\parskip}{0.45em}}
  {\par\medskip\hrule\medskip}

% --- Mots-clés ---------------------------------------------------------------
%
% En anglais, en clôture du résumé de la lecture modernisée : le vocabulaire
% moderne sous lequel chercher ce que les feuillets construisent. Le manifeste
% du site (`npm run manifest`) les extrait pour étiqueter la cote entière —
% cette ligne est la source unique des tags, il n'y a pas de fichier à côté.
\newcommand{\keywords}[1]{\par\smallskip\noindent
  {\footnotesize\textsc{Keywords} --- \textit{#1}}\par}

% --- L'appareil critique ----------------------------------------------------

% Le numéro de feuillet, en marge. C'est l'ancre qui permet de régler un
% désaccord sur une formule en regardant *un* feuillet plutôt qu'une chemise
% de six cents. Le site s'en sert aussi pour tourner les pages du fac-similé
% quand on fait défiler la transcription.
\newcommand{\leaf}[1]{%
  \marginnote{\footnotesize\color{gray!70}\textsf{#1}}%
  \label{leaf:#1}\ignorespaces}

% Illisible : on ne devine pas. Un mot inventé qui se lit comme les autres est
% le pire résultat possible.
\newcommand{\ill}{\textcolor{red!60!black}{[\,\ldots\,]}}

% Lecture douteuse : le mot est donné, et signalé comme incertain.
\newcommand{\uncertain}[1]{\uline{#1}}

% Ajout de l'éditeur — une lettre manquante, un mot sous-entendu.
\newcommand{\add}[1]{\textcolor{blue!50!black}{[#1]}}

% Biffé par Grothendieck lui-même. Ce qu'il a rayé fait partie du document :
% une rature dit souvent où la pensée a changé de direction.
\newcommand{\struck}[1]{\sout{#1}}

% Note du transcripteur, sur l'état matériel du feuillet.
\newcommand{\note}[1]{\par\smallskip{\small\color{gray!60!black}\textit{[#1]}}\par\smallskip}

% Note marginale de Grothendieck, distincte du corps du texte.
\newcommand{\marginal}[1]{\par\smallskip\noindent\hspace{1em}%
  {\small\color{gray!60!black}#1}\par\smallskip}

% --- Titre ------------------------------------------------------------------

\AtBeginDocument{%
  \begin{center}
    {\large\bfseries Fonds Alexandre Grothendieck --- cote n\textsuperscript{o}~\grothFolder}\\[2pt]
    {\itshape\grothFoldertitle}\\[4pt]
    {\small feuillets \grothFirst--\grothLast\ (lot \grothBatch)}\\[2pt]
    {\footnotesize\color{gray!60!black}Transcription établie d'après le fac-similé numérisé
     par l'Université de Montpellier.\\
     Les lectures incertaines sont soulignées, les illisibles notées [\,\ldots\,],
     les ajouts entre crochets.}
  \end{center}
  \vspace{1em}}

\endinput


\folder{108}
\pages{1}{5}
\dating{[à partir de 1973]}
\watermark{Édition de démonstration\\interprétation personnelle de l'œuvre}
\foldertitle{Divers, liste questions : notes manuscrites (s.d.) — lecture modernisée du dossier entier}

\begin{document}

\section*{Résumé}

\begin{resume}
Une chemise portant deux mots au crayon --- « divers, liste questions » --- et,
dedans, dix-sept questions numérotées. Rien d'autre : pas une démonstration,
pas un énoncé affirmatif, pas même une phrase complète dans plusieurs cas. Un
inventaire de ce qui reste à faire.

Ce qui reste à faire, ici, concerne un projet dont il faut dire d'abord ce
qu'il est. Depuis Poincaré, on sait attacher à un espace des invariants
algébriques --- des groupes qui comptent ses trous --- et on sait que deux
espaces très différents d'aspect peuvent avoir exactement les mêmes. On dit
alors qu'ils ont le même \emph{type d'homotopie}. La question naturelle est
alors : de quoi le type d'homotopie est-il vraiment fait ? Si des objets aussi
dissemblables qu'une sphère de caoutchouc et un modèle combinatoire fait de
triangles recollés portent la même information, c'est que l'information ne
tient ni au caoutchouc ni aux triangles. À quoi tient-elle ?

La réponse que ces pages poursuivent est catégorique. On ne cherche pas un
meilleur espace : on cherche à caractériser les \emph{catégories} qui savent
porter les types d'homotopie. Une telle catégorie, munie de la classe de
flèches qu'on décide d'y tenir pour des équivalences, est ce que Grothendieck
appelle un \emph{modélisateur}. Les triangles en donnent un ; les espaces
topologiques en donnent un autre ; et la découverte qui organise tout le projet
est que la catégorie des \emph{petites catégories} en donne un aussi. Une
catégorie, objet a priori sans aucune géométrie, suffit à porter un type
d'homotopie.

Une fois cela vu, la question devient : lesquelles ? Quelles petites catégories
$A$ ont la propriété que les préfaisceaux sur $A$ modélisent l'homotopie ? Ce
sont les \emph{catégories test}, et les questions de cette liste tournent
presque toutes autour d'elles : est-ce que $\Delta^{\circ}$ en est une, quelles
propriétés d'exactitude a la catégorie homotopique, à quoi ressemble une
fibration dans ce cadre, que devient la construction de Dold--Puppe.

La liste est écrite en deux langues --- les dix premières questions en anglais,
les sept dernières en français --- et sans transition entre les deux. Elle a
été écrite d'un trait, puis reprise : deux numéros sont intervertis par une
flèche, plusieurs lignes sont ajoutées entre les lignes. C'est un état
intermédiaire d'un travail, pris tel quel.

\keywords{modelizer, test category, weak equivalence, contractor, Dold-Puppe, cointegration, homotopy type of a topos}
\end{resume}

\pagerange{2}{2}
\section*{Le vocabulaire de la liste}

Quatre notations reviennent partout et n'ont nulle part de définition sur la
page. On les fixe ici une fois pour toutes ; elles valent pour tout ce qui
suit.

\begin{itemize}
\item $A$ désigne une petite catégorie, et $\hat{A}$ la catégorie des
préfaisceaux sur $A$, c'est-à-dire des foncteurs $A^{\circ} \to
(\mathrm{Ens})$.
\item $(\mathrm{Cat})$ est la catégorie des petites catégories.
\item $\underline{W}$ est une classe de flèches déclarées équivalences faibles ;
$W_{A}$ est celle de $\hat{A}$, et $(\mathrm{Hot}_{\underline{W}})$ la
catégorie obtenue en inversant $\underline{W}$.
\item $\Delta$ est la catégorie des simplexes, $\hat{\Delta}$ celle des
ensembles simpliciaux.
\end{itemize}

Un \emph{modélisateur} est un couple $(M, \underline{W})$ tel que le localisé
$\underline{W}^{-1}M$ soit équivalent à la catégorie homotopique des espaces. La
liste emploie le mot sans le définir, et distingue en passant les modélisateurs
\emph{élémentaires} --- ceux de la forme $\hat{A}$ --- des autres.\footnote{La
question 10 emploie aussi « strict », entre parenthèses et suivi d'un point
d'interrogation, sans le définir davantage. On ne lui donne pas de sens ici.}

\pagerange{2}{2}
\section*{Les dix premières questions}

\subsection*{Exactitude et structure de modèles (1 à 5)}

Les questions 1 et 2 portent sur les conditions d'asphéricité de type Čech et
sur la saturation des classes $\underline{W}$ et $W_{A}$.\footnote{Une classe
d'équivalences faibles est saturée si toute flèche devenue inversible dans le
localisé y était déjà. C'est la condition minimale pour que la classe soit
déterminée par le localisé qu'elle produit ; la liste la pose comme une
question, ce qui suppose qu'elle n'allait pas de soi.}

La question 3 demande les propriétés homotopiques de $(\mathrm{Cat})$ et les
suites de Dold--Puppe, puis, sur une ligne à part :

\begin{quote}
$(\mathrm{Cat})$ est-elle une catégorie de modèles ?
\end{quote}

La question 5 pose la même chose pour $\hat{A}$ : quand $\hat{A}$ est-elle une
catégorie de modèles (fermée) ?\footnote{Les numéros 4 et 5 sont intervertis sur
la page --- le 5 est écrit d'abord --- et une flèche tracée dans la marge
rétablit l'ordre. On suit la flèche. Le mot « fermée » traduit \emph{closed} au
sens de Quillen : une structure de modèles dont deux des trois classes
déterminent la troisième.} Et la question 4 demande les propriétés
d'exactitude de $(\mathrm{Hot}_{\underline{W}})$ ainsi que celles des deux
foncteurs canoniques
\[
(\mathrm{Cat}) \longrightarrow (\mathrm{Hot}_{\underline{W}}),
\qquad
\hat{A} \longrightarrow (\mathrm{Hot}_{\underline{W}}).
\]

Ces trois questions sont une seule, posée trois fois : jusqu'où la machinerie de
Quillen s'installe-t-elle sur les catégories qui nous intéressent ? Poser la
question pour $(\mathrm{Cat})$ suppose déjà acquis que $(\mathrm{Cat})$ est un
modélisateur, ce que la liste ne redémontre pas.

\subsection*{Une réponse, la seule de la liste (6)}

La question 6 est la seule qui reçoive sa réponse sur la page même :

\begin{quote}
La catégorie $\Delta^{\circ}$ est-elle une catégorie test, ou non ? \emph{non}
\end{quote}

Le mot est souligné.\footnote{La liste ne dit pas pourquoi. On ne comble pas :
un contre-exemple ou un argument serait le nôtre, et cette édition n'en
fabrique pas.} La question se poursuit par une liste de catégories cubiques et
simpliciales voisines --- $\Delta_{f}$, et trois variantes d'un même signe que
la page ne nomme pas ---, puis par un alinéa qui dit seulement :
\emph{cubical games}.\footnote{Le signe en question est un ovale épais tracé
d'un seul trait. L'alinéa suivant parle de jeux cubiques, ce qui invite à y lire
la catégorie cubique $\square$, et c'est ainsi qu'on l'a rendu dans la
transcription ; mais la page ne le nomme pas et l'identification reste la nôtre.
Les trois variantes portent respectivement un exposant, un tilde et un indice
$f$, sans qu'on sache ce qu'ils distinguent.}

\subsection*{Invariants, intégration, réalisation (7 à 9)}

La question 7 demande des invariants canoniques attachés à un modélisateur
élémentaire $\hat{A}$, puis : Dold--Puppe axiomatisé ? et le cas
cubique.\footnote{Un mot barré précède « invariants » et se lit
vraisemblablement \emph{cohomology} ; la rature est nette et on ne le compte
pas dans l'énoncé.}

La question 8 est la plus riche des dix. Elle demande une notion
d'\emph{intégrale} de types d'homotopie --- une limite projective, écrit-elle,
avec la flèche vers la gauche --- et sa notion duale ; puis, sur la même ligne,
la structure d'une « catégorie d'homotopie » : la structure interne de
$(\mathrm{Hot})$ et de ce qui lui ressemble, et une \emph{structure de
contractibilité}.\footnote{Ces trois derniers mots sont ajoutés entre les
lignes, d'un trait plus fin. C'est le seul endroit de la liste où figure la
notion de contractibilité, qui reparaîtra pourtant en français à la question 16
sous le nom de contracteur.}

La question 9 demande si les espaces topologiques forment un modélisateur ---
le mot est souligné sur la page ---, quels foncteurs test prennent leurs
valeurs dans les espaces, et ce qu'est la réalisation géométrique dans ce
cadre. La question est moins naïve qu'elle en a l'air : dans le programme des
modélisateurs, les espaces topologiques ne sont plus le point de départ mais un
cas parmi d'autres, dont il faut vérifier qu'il entre bien dans le cadre.

\subsection*{Types pointés (10)}

La question 10 demande ce que deviennent les objets pointés et les objets en
groupes : si $(M, \underline{W})$ est un modélisateur, la catégorie
$\mathrm{Gr}(M)$ de ses objets en groupes est-elle un modélisateur pour les
types d'homotopie pointés 0-connexes ? Et la liste ajoute : lié à
Dold--Puppe ?\footnote{La fin de la question court jusqu'au bord de la feuille
et une rature en recouvre une partie ; la transcription le signale. Le rapport
à Dold--Puppe est en effet celui qui rend la question naturelle --- la
correspondance de Dold--Kan identifie objets simpliciaux abéliens et complexes
de chaînes --- mais la page ne développe pas et l'on n'en dit pas plus.}

\pagerange{3}{4}
\section*{Les sept questions françaises}

Sans transition ni annonce, la liste passe au français à la question 11.

\subsection*{Fibrations (11 et 12)}

La question 11 cherche une caractérisation des fibrations dans
$(\mathrm{Cat})$. La définition visée est claire : un morphisme $X' \to X$ tel
que le changement de base le long de lui transforme les équivalences faibles en
équivalences faibles. La question est de savoir si cette condition, lourde à
vérifier telle quelle, se teste sur une famille restreinte --- les inclusions
$e \to Y$ d'un objet final, ou bien les composées $\Delta_{0} \to \Delta_{n}
\to X$.\footnote{Une abréviation de la page se lit \emph{il f. et s.}, c'est-à-dire
« il faut et il suffit », et la transcription la donne comme incertaine. La
lecture retenue est celle qui rend la phrase grammaticale et l'argument
cohérent ; une autre lecture changerait la force de l'énoncé, pas son objet.}

La question 12 demande comment interpréter, dans l'axiomatique de Quillen, la
condition qui vient d'être posée. Elle l'écrit sous forme de carré cartésien :

\begin{tikzcd}[column sep=small, row sep=small, nodes={font=\scriptsize}]
X' \arrow[r, "g'"] \arrow[d, "f'"'] & X \arrow[d, "f"]\\
Y' \arrow[r, "g"] & Y
\end{tikzcd}

et la condition est : si $f$ est une fibration et $g \in \underline{W}$, alors
$g' \in \underline{W}$. C'est la stabilité des équivalences faibles par
changement de base le long d'une fibration --- ce qu'on appelle aujourd'hui la
propreté à droite d'une structure de modèles.\footnote{Ce nom est le nôtre. La
liste écrit la condition et demande son interprétation ; elle ne la baptise pas,
et rien n'indique qu'elle ait eu ce nom à ses yeux.}

\subsection*{Asphéricité et 0-connexité (13 et 14)}

La question 13 demande si une catégorie asphérique dont le topos est
0-connexe l'est \emph{totalement}, c'est-à-dire si la 0-connexité est stable par
produit : le produit de deux objets 0-connexes est-il 0-connexe ? C'est le
même passage du local au global que la question 11 posait pour les fibrations,
et c'est le mécanisme qui distingue partout dans ce programme l'asphérique du
totalement asphérique.\footnote{C'est la question la plus surchargée de la
liste : trois mots au moins y sont écrits l'un sur l'autre au début, un
\emph{test} est barré au profit d'un \emph{aspherical}, et la parenthèse finale
est reprise deux fois. La transcription donne ce qui se lit. L'énoncé restitué
ici est donc plus net que la page, et c'est un choix : une question posée trois
fois de trois manières se laisse résumer, mais le résumé est nôtre.}

La question 14 tient en une ligne : les $\mathrm{Hom}$ dans $(\mathrm{Hot})$ et
les variantes diagrammatiques.\footnote{Le premier mot est souligné et se lit
\emph{Hom} sans certitude absolue ; la transcription le marque.}

\subsection*{Cointégration (15)}

La question 15 demande une théorie de la \emph{cointégration} sur un topos, et
d'intégration pour certains morphismes de topos, avec une suite spectrale de la
forme
\[
\underline{\Pi}_{*}\Bigl(R\!\int_{*}(\xi)\Bigr) \longleftarrow
R^{i}f_{*}\bigl(\underline{\Pi}_{j}(\mathcal{G})\bigr)
\]
pour un type d'homotopie relatif sur un topos $X$.\footnote{La formule est
écrite très vite et plusieurs de ses indices sont ambigus : l'écriture ne
distingue pas nettement le $i$ du $j$, et le rapport entre l'étoile du membre
de gauche et les deux indices de droite n'est pas explicité. On la reproduit
telle qu'elle se lit, sans harmoniser les indices, parce qu'une suite spectrale
dont on aurait « corrigé » les degrés ne serait plus la sienne. Elle est à lire
comme une forme visée, non comme un énoncé démontré.} C'est la question 8 ---
l'intégrale de types d'homotopie --- reprise en français et transportée sur les
topos, avec cette fois un candidat d'énoncé.

\subsection*{Contracteurs (16)}

La question 16 est la plus développée de la liste, et la seule qui se
subdivise. Soit $A$ un \emph{contracteur} non trivial, et soient les deux
foncteurs

\begin{tikzcd}[column sep=large, nodes={font=\scriptsize}]
\hat{A} \arrow[r, bend left=20, "i_{A}"] & (\mathrm{Cat}) \arrow[l, bend left=20, "j_{A}"]
\end{tikzcd}

Quatre questions s'enchaînent.

\begin{enumerate}
\item Ces deux foncteurs induisent-ils des équivalences quasi-inverses l'une de
l'autre entre les localisés par les \emph{homotopismes} ?
\item Est-il vrai qu'un $X$ de $\hat{A}$ est contractile si et seulement si
$A_{/X}$ l'est ?\footnote{Une abréviation suit « cont. » et se lit
\emph{t.-l.}, sans qu'on sache la développer ; la transcription la marque
incertaine. L'énoncé donné ici est celui que le reste de la phrase impose.}
\item Est-il vrai qu'une petite catégorie $C$ est contractile si et seulement si
$A_{/C} = j_{A}(C)$ l'est ?
\item Est-il vrai que pour $X, Y$ dans $\hat{A}$, la flèche canonique
$i_{A}(X \times Y) \to i_{A}(X) \times i_{A}(Y)$ est un homotopisme ?
\end{enumerate}

Puis, souligné : \emph{variantes} commutatives de ces questions.

Ces quatre questions ont un point commun, et c'est lui qui fait leur unité : ce
sont toutes des questions de \emph{conservation}. On demande si la
contractilité, la contractilité encore, et le produit, se transportent
intacts d'un côté à l'autre de l'adjonction. Ce que $i_{A}$ et $j_{A}$
préservent est exactement ce qui autorisera à travailler indifféremment dans
$\hat{A}$ ou dans $(\mathrm{Cat})$.

\subsection*{Et pour finir (17)}

Cinq mots : Dold--Puppe, Whitehead, etc.

\pagerange{5}{5}
\section*{Le feuillet mobile}

Un feuillet déchiré, perforé, écrit plus vite et raturé sur trois lignes
pleines, ferme le dossier. Il porte un fragment d'argument dont le squelette est
lisible et le détail non.\footnote{La transcription y compte une dizaine de
passages illisibles sur une quinzaine de lignes. Ce qui suit ne restitue que ce
qui se lit ; on ne comble aucun des blancs, et l'on n'en tire aucune
conséquence.}

Ce qui s'en dégage est ceci. Un topos strictement totalement asphérique n'est
pas modélisant, sauf s'il provient d'un $\hat{A}$. La question qui suit --- avec
un \emph{Cat-test} souligné et suivi d'un point d'exclamation --- est de savoir
s'il suffit de trouver une catégorie $A$ à produits finis vérifiant une
condition d'unicité sur son objet initial. Et la dernière ligne, qui est la
seule que la rature épargne entièrement, conclut :

\begin{quote}
Mais si $A$ est un contracteur, alors il est [\ldots] si et seulement si
$A \cong \Delta_{0}$.
\end{quote}

Un mot manque entre « il est » et la condition.\footnote{Le mot illisible
qualifie $A$ et commence par ce qui pourrait être « cat. » suivi d'un
adjectif ; on ne le devine pas. $\Delta_{0}$ est la catégorie ponctuelle, de
sorte que l'énoncé dit, quelle que soit la propriété manquante, que seul le
contracteur trivial la possède --- ce qui est bien un énoncé de rigidité, et
qui répond à la question 16 dans le sens négatif pour le cas dégénéré. Cette
lecture est la nôtre.}

\end{document}
